% Part: model-theory
% Chapter: interpolation
% Section: interpolation-proof

\documentclass[../../../include/open-logic-section]{subfiles}

\begin{document}

\olfileid[hi]{mod}{int}{prf}

\olsection{क्रेग का अंतर्वेशन प्रमेय}

\begin{thm}[क्रेग का अंतर्वेशन प्रमेय]
\ollabel{thm:interpol} 
यदि $\Entails !A \lif !B$, तो ऐसा !!a{sentence} $!C$ है कि
$\Entails !A \lif !C$ और $\Entails !C \lif !B$, और प्रत्येक
!!{constant}, !!{function} और !!{predicate} ($\eq$ के अतिरिक्त), जो
$!C$ में आता है, $!A$ और~$!B$ दोनों में आता है। !!{sentence} $!C$ को
$!A$ और~$!B$ का
\emph{अंतर्वेशक} कहते हैं।
\end{thm}

\begin{proof}
मान लें कि $\Lang{L}_1$, $!A$ की भाषा है और $\Lang{L}_2$, $!B$ की भाषा
है। $\Lang{L}_0 = \Lang{L}_1 \cap \Lang{L}_2$ मानें। प्रत्येक
$i \in \{0, 1, 2 \}$ के लिए, $\Lang{L}'_i$ को $\Lang{L}_i$ में अनंततः
अनेक नए !!{constant}s $\Obj c_0,
\Obj c_1, \Obj c_2, \dots$ जोड़कर
प्राप्त करें।

यदि $!A$ असंतुष्टनीय है, तो $\lexists[x][\eq/[x][x]]$ एक अंतर्वेशक है।
यदि $\lnot !B$ असंतुष्टनीय है (और इसलिए $!B$ वैध है), तो
$\lexists[x][\eq[x][x]]$ एक अंतर्वेशक है। अतः हम यह भी मान सकते हैं कि
$!A$ और $\lnot !B$ दोनों संतुष्टनीय हैं।

अंतर्वेशन प्रमेय का प्रतिपक्ष सिद्ध करने के लिए मान लें कि $!A$ और $!B$
का कोई अंतर्वेशक नहीं है। दूसरे शब्दों में, मान लें कि $\{!A\}$ और
$\{\lnot !B\}$, $\Lang{L}_0$ में अपृथक्करणीय हैं।

हमारा लक्ष्य युग्म $(\{ !A \}, \{\lnot!B\})$ को किसी अधिकतम
अपृथक्करणीय युग्म $(\Gamma^*, \Delta^*)$ तक विस्तृत करना है। $!A_0$,
$!A_1$, $!A_2$, \dots को $\Lang{L}_1$ के !!{sentence}s की गणना और
$!B_0$, $!B_1$, $!B_2$, \dots को $\Lang{L}_2$ के !!{sentence}s की
गणना मानें। हम !!{sentence}s के समुच्चयों के दो वर्धमान अनुक्रम
$(\Gamma_n, \Delta_n)$, जहाँ $n \ge 0$, निम्न प्रकार परिभाषित करते हैं।
$\Gamma_0 = \{ !A\}$ और $\Delta_0 = \{\lnot !B \}$ रखें। मान लें कि
$(\Gamma_n, \Delta_n)$ पहले से परिभाषित हैं; तब $\Gamma_{n+1}$ और
$\Delta_{n+1}$ को इस प्रकार परिभाषित करें:
\begin{enumerate}
\item यदि $\Gamma_n \cup \{!A_n \}$ और $\Delta_n$, $\Lang{L}'_0$ में
  अपृथक्करणीय हैं, तो $!A_n$ को $\Gamma_{n+1}$ में रखें। इसके अतिरिक्त,
  यदि $!A_n$ कोई अस्तित्ववाची !!{formula} $\lexists[x][!S]$ है, तो ऐसा
  नया !!{constant} $c$ चुनें जो $\Gamma_n$, $\Delta_n$, $!A_n$ या
  $!B_n$ में नहीं आता, और $\Subst{!S}{c}{x}$ को $\Gamma_{n+1}$ में रखें।
\item यदि $\Gamma_{n+1}$ और $\Delta_n \cup \{!B_n \}$,
  $\Lang{L}'_0$ में अपृथक्करणीय हैं, तो $!B_n$ को $\Delta_{n+1}$ में
  रखें। इसके अतिरिक्त, यदि $!B_n$ कोई अस्तित्ववाची !!{formula}
  $\lexists[x][!S]$ है, तो ऐसा नया !!{constant} $c$ चुनें जो
  $\Gamma_{n+1}$, $\Delta_n$, $!A_n$ या $!B_n$ में नहीं आता, और
  $\Subst{!S}{c}{x}$ को $\Delta_{n+1}$ में रखें।
\end{enumerate}
अंत में, परिभाषित करें:
\begin{align*}
  \Gamma^* & = \bigcup_{n\ge 0} \Gamma_n, & 
  \Delta^* & = \bigcup_{n\ge 0} \Delta_n.
\end{align*}
अब $n$ पर समकालिक आगमन द्वारा हम सिद्ध कर सकते हैं:
\begin{enumerate}
\item\ollabel{part-a} $\Gamma_n$ और $\Delta_n$, $\Lang{L}'_0$ में
  अपृथक्करणीय हैं;
\item\ollabel{part-b} $\Gamma_{n+1}$ और $\Delta_n$, $\Lang{L}'_0$ में
    अपृथक्करणीय हैं।
\end{enumerate}
\olref{part-a} का आधार \olref[sep]{lem:sep1} से मिलता है। भाग
\olref{part-b} के लिए हमें तीन स्थितियाँ अलग करनी होंगी:
\begin{enumerate}
\item यदि $\Gamma_0 \cup \{!A_0 \}$ और $\Delta_0$ पृथक्करणीय हैं, तो
  $\Gamma_1 = \Gamma_0$ और \olref{part-b}, \olref{part-a} ही है;
\item यदि $\Gamma_1 = \Gamma_0 \cup\{ !A_0\}$, तो रचना से $\Gamma_1$
  और $\Delta_0$ अपृथक्करणीय हैं।
\item वह स्थिति शेष है जिसमें $!A_0$ अस्तित्ववाची है, अतः
  $\Gamma_1 = \Gamma_0 \cup \{ \lexists[x][!S], \Subst{!S}{c}{x}
  \}$। रचना से $\Gamma_0 \cup \{ \lexists[x][!S]\}$ और $\Delta_0$
  अपृथक्करणीय हैं, इसलिए \olref[sep]{lem:sep2} द्वारा
  $\Gamma_0 \cup \{ \lexists[x][!S], \Subst{!S}{c}{x} \}$ और
  $\Delta_0$ भी अपृथक्करणीय हैं।
\end{enumerate}
इससे ऊपर के \olref{part-a} और \olref{part-b} के लिए आगमन का आधार पूरा
होता है। अब आगमन-चरण लें। \olref{part-a} के लिए, यदि
$\Delta_{n+1} = \Delta_n \cup \{ !B_n \}$, तो रचना से
$\Gamma_{n+1}$ और $\Delta_{n+1}$ अपृथक्करणीय हैं ($!B_n$ के
अस्तित्ववाची होने पर भी, \olref[sep]{lem:sep2} द्वारा); यदि
$\Delta_{n+1} = \Delta_n$ (क्योंकि $\Gamma_{n+1}$ और
$\Delta_n \cup \{!B_n\}$ पृथक्करणीय हैं), तो हम \olref{part-b} पर
आगमन-परिकल्पना का उपयोग करते हैं। \olref{part-b} के आगमन-चरण के लिए, यदि
$\Gamma_{n+2} = \Gamma_{n+1} \cup
\{!A_{n+1} \}$, तो रचना से
$\Gamma_{n+2}$ और $\Delta_{n+1}$ अपृथक्करणीय हैं ($!A_{n+1}$ के
अस्तित्ववाची होने पर भी, \olref[sep]{lem:sep2} द्वारा); और यदि
$\Gamma_{n+2} = \Gamma_{n+1}$, तो हम अभी सिद्ध किए गए
\olref{part-a} के आगमन-विषयक मामले का उपयोग करते हैं। इससे
\olref{part-a} और \olref{part-b} पर आगमन पूरा होता है।

इससे $\Gamma^*$ और $\Delta^*$ के अपृथक्करणीय होने का निष्कर्ष निकलता
है; अन्यथा संहतता से कोई $n \ge 0$ होता जो $\Gamma_n$ और $\Delta_n$
को पृथक करता, जो \olref{part-a} के विरुद्ध है। विशेषतः, $\Gamma^*$ और
$\Delta^*$ संगत हैं: क्योंकि यदि पहला या दूसरा असंगत होता, तो उन्हें
क्रमशः $\lexists[x][\eq/[x][x]]$ या $\lforall[x][\eq[x][x]]$ पृथक
करते।

अब हम दिखाते हैं कि $\Gamma^*$, $\Lang{L}'_1$ में अधिकतम संगत है और इसी
प्रकार $\Delta^*$, $\Lang{L}'_2$ में अधिकतम संगत है। पहले के लिए मान लें
कि $!A_n \notin \Gamma^*$ और $\lnot !A_n
\notin \Gamma^*$, किसी $n \ge 0$ के लिए। यदि $!A_n \notin \Gamma^*$, तो
$\Gamma_n \cup \{!A_n \}$, $\Delta_n$ से पृथक्करणीय है, और इसलिए ऐसा
$!C \in \Lang{L}'_0$ है कि दोनों:
\begin{align*}
  \Gamma^* & \Entails !A_n \lif !C, & 
  \Delta^* & \Entails \lnot !C.
\end{align*}
इसी प्रकार, यदि $\lnot !A_n \notin \Gamma^*$, तो ऐसा
$!C' \in
\Lang{L}'_0$ है कि दोनों:
\begin{align*}
  \Gamma^* & \Entails \lnot !A_n \lif !C', & 
  \Delta^* & \Entails \lnot !C'.
\end{align*}
प्रतिज्ञप्तिक तर्क से $\Gamma^* \Entails !C \lor !C'$ और
$\Delta^* \Entails \lnot (!C \lor !C')$, इसलिए $!C \lor
!C'$, $\Gamma^*$ और $\Delta^*$ को पृथक करता है। ऐसा ही तर्क स्थापित
करता है कि $\Delta^*$ अधिकतम है।

अंततः हम दिखाते हैं कि $\Gamma^* \cap \Delta^*$, $\Lang{L}'_0$ में
अधिकतम संगत है। वह स्पष्टतः संगत है, क्योंकि वह संगत समुच्चयों का
सर्वनिष्ठ है। अधिकतमता दिखाने के लिए $!S \in
\Lang{L}'_0$ लें। अब $\Gamma^*$, $\Lang{L'_1}
\supseteq \Lang{L'_0}$ में अधिकतम है और इसी
प्रकार $\Delta^*$, $\Lang{L'_2} \supseteq \Lang{L'_0}$ में अधिकतम है।
इससे या तो $!S \in \Gamma^*$ या $\lnot !S \in \Gamma^*$, और या तो
$!S \in \Delta^*$ या $\lnot !S \in \Delta^*$। यदि $!S \in
\Gamma^*$ और $\lnot !S \in \Delta^*$, तो $!S$, $\Gamma^*$ और
$\Delta^*$ को पृथक करता; और यदि $\lnot !S \in
\Gamma^*$ और $!S \in \Delta^*$, तो $\Gamma^*$ और $\Delta^*$ को
$\lnot !S$ पृथक करता। अतः या तो $!S \in
\Gamma^* \cap \Delta^*$ या $\lnot !S \in \Gamma^* \cap \Delta^*$,
और $\Gamma^* \cap \Delta^*$ अधिकतम है।

चूँकि $\Gamma^*$ अधिकतम संगत है, उसका ऐसा प्रतिरूप $\Struct{M}'_1$ है
जिसका !!{domain} $\Domain{M'_1}$ ठीक उन अवयवों $\Assign{c}{M'_1}$ से
बना है जो !!{constant}s की व्याख्या करते हैं---ठीक पूर्णता प्रमेय के
प्रमाण की तरह (\olref[fol][com][cth]{thm:completeness})। इसी प्रकार,
$\Delta^*$ का ऐसा प्रतिरूप $\Struct{M}'_2$ है जिसका !!{domain}
$\Domain{M'_2}$, !!{constant}s की व्याख्याओं $\Assign{c}{M'_2}$ से
दिया जाता है।

$\Struct{M_1}$ को $\Struct{M'_1}$ से $\Lang{L'_1} \setminus \Lang{L'_0}$
के !!{constant}s, !!{function}s और !!{predicate}s की
व्याख्याएँ हटाकर प्राप्त करें, और $\Struct{M_2}$ के लिए भी ऐसा ही करें।
तब प्रतिचित्र $h \colon M_1 \to M_2$, जिसे
$h(\Assign{c}{M'_1}) = \Assign{c}{M'_2}$ द्वारा परिभाषित किया गया है,
$\Lang{L}'_0$ में एक समरूपता है, क्योंकि जैसा
दिखाया गया है, $\Gamma^* \cap \Delta^*$, $\Lang{L}'_0$ में अधिकतम संगत
है। यह इसलिए निकलता है कि कोई भी $\Lang{L}'_0$-!!{sentence} या तो
$\Gamma^*$ और $\Delta^*$ दोनों में है या दोनों में नहीं है: इसलिए
$\Assign{c}{M'_1} \in
\Assign{P}{M'_1}$ तभी और केवल तभी जब
$\Atom{P}{c} \in \Gamma^*$, तभी और केवल तभी जब
$\Atom{P}{c} \in \Delta^*$, और तभी और केवल तभी जब
$\Assign{c}{M'_2} \in
\Assign{P}{M'_2}$। समरूपताओं द्वारा संतुष्ट की
जाने वाली अन्य शर्तें भी इसी प्रकार स्थापित की जा सकती हैं।

अब एक प्रतिरूप $\Struct{M}$ को !!{language}
$\Lang{L_1} \cup \Lang{L_2}$ के लिए निम्न प्रकार परिभाषित करें:
\begin{enumerate}
\item !!{domain} $\Domain{M}$ ठीक $\Domain{M_2}$ है, अर्थात् सभी अवयवों
  $\Assign{c}{M'_2}$ का समुच्चय;
\item यदि !!a{predicate}~$P$, $\Lang{L_2} \setminus
  \Lang{L_1}$ में
  है, तो $\Assign{P}{M} = \Assign{P}{M'_2}$;
\item यदि कोई विधेय $P$, $\Lang{L}_1\setminus \Lang{L}_2$ में है, तो
  $\Assign{P}{M} = h(\Assign{P}{M'_2})$, अर्थात्
  $\tuple{\Assign{c_1}{M'_2}, \dots, \Assign{c_n}{M'_2}} \in
  \Assign{P}{M}$ तभी और केवल तभी जब
  $\tuple{\Assign{c_1}{M'_1}, \dots,
  \Assign{c_n}{M'_1}} \in \Assign{P}{M'_1}$।
\item यदि !!a{predicate} $P$, $\Lang{L}_0$ में है, तो
  $\Assign{P}{M} =
  \Assign{P}{M'_2} = h(\Assign{P}{M'_1})$।
\item $\Lang{L}_1 \cup \Lang{L}_2$ के !!^{function}s, जिनमें
  !!{constant}s भी सम्मिलित हैं, इसी प्रकार सँभाले जाते हैं।
\end{enumerate}

अंततः !!{formula}s पर आगमन द्वारा दिखाया जाता है कि $\Struct{M}$,
$\Struct{M'_1}$ से $\Lang{L'_1}$ के सभी !!{formula}s पर और
$\Struct{M'_2}$ से $\Lang{L'_2}$ के सभी !!{formula}s पर सहमत है।
विशेषतः, $\Struct{M} \Entails \Gamma^* \cup \Delta^*$, जिससे
$\Struct{M}
\Entails !A$ और $\Struct{M} \Entails \lnot!B$, और
$\not\Entails !A \lif !B$। इससे क्रेग के अंतर्वेशन प्रमेय का प्रमाण
पूरा होता है।
\end{proof}

\end{document}
